\subsection{簿册与粮船：开元秩序怎样进入乡里和河道}
721年宇文融主持\eventref{ST-721-YUWEN-RONG}{括户}，朝廷希望找回逃户、隐漏户和田亩。新增数字是行政口径，不能当作真实人口或财富的照片；“隐户”也不必然是有意欺瞒，有人因迁徙、婚姻、战乱或依附关系离开旧籍。可是一张重填的簿册会改变一户人的位置：里正获得核查权，家庭可能重新面对租庸调，原来提供保护的亲族或寺院也要重新计算责任。国家在这里不是突然看清乡村，而是在文书、田界和逃避之间制造新的摩擦。

733年裴耀卿整顿\eventref{ST-733-PEI-YAOQING}{江淮转运}。传记所称漕粮数额与效率不能直接外推；一粒粮从江淮到关中要经过装船、验收、换运、入仓和再发运。水位、堤岸、船工和沿路治安都可能改变它的行程。对长安而言，这是一条供给线；对江淮的船户、盐商与州县而言，它意味着征用、市场机会和风险同时增加。开元的秩序因此不是中心单方面得到粮食，而是地方不断承担把粮食变成财政的劳动。

天宝六载征南诏失利、十载怛罗斯受挫，\eventref{ST-747-NANZHAO}{云南山路}与\eventref{ST-751-TALAS}{中亚战场}以不同方式耗费唐的资源。同年安禄山兼领三镇，\eventref{ST-751-AN-LUSHAN-COMMANDS}{东北军镇资源集中}。唐方战史不足以精确统计伤亡，外部材料也不能填平所有空白；但道路、马匹和军粮已表明，朝廷不能无限延长每一条边线。后来危机的压力，不是从一位将领的性格忽然生出，而是在这些簿册、粮船和边地补给中慢慢堆积。
